% ---------------------------------------------------------------
%  Muc II.4 -- Huong dan su dung + loi thuong gap
% ---------------------------------------------------------------
\subsection{Hướng dẫn sử dụng hệ thống}

Mục này viết cho người dùng thật --- giáo viên và học sinh --- nên trình bày
theo từng thao tác, mỗi thao tác kèm một hình chụp màn hình thực tế.

\subsubsection{Địa chỉ truy cập}

\begin{center}
\setlength{\fboxsep}{10pt}
\fbox{\begin{minipage}{0.7\linewidth}
\centering
\textbf{\large \url{https://nganhangdechv.tech}}\\[8pt]
\qrcode[height=3cm]{https://nganhangdechv.tech}\\[6pt]
{\small Quét mã để vào hệ thống}
\end{minipage}}
\end{center}

Hệ thống chạy trực tiếp trên trình duyệt, dùng được cả trên máy tính và điện
thoại, không phải cài đặt phần mềm gì thêm. Mọi mã QR trong báo cáo này đều do
\LaTeX\ sinh ra từ chính đường dẫn in bên cạnh, nên không có nguy cơ mã và
đường dẫn lệch nhau.

\subsubsection{Hướng dẫn dành cho học sinh}

\textbf{Bước 1 --- Tạo tài khoản và đăng nhập.} Vào trang chủ, chọn
\textit{Đăng ký}, điền họ tên, email và mật khẩu; hoặc bấm \textit{Đăng nhập
bằng Google} để dùng luôn tài khoản Gmail của mình. Học sinh đã có tài khoản
thì đăng nhập thẳng. Nếu quên mật khẩu, dùng chức năng \textit{Quên mật khẩu}
để nhận thư đặt lại.

\chenanh[0.65]{09-dang-ky-dang-nhap.png}{Màn hình đăng ký và đăng nhập của học sinh}

\textbf{Bước 2 --- Chọn lớp.} Lần đầu đăng nhập, học sinh chọn lớp mình đang
học trong danh sách. Việc này để hệ thống biết phạm vi chương trình và để giáo
viên nhận đúng kết quả của lớp mình. Có thể quay lại đổi lớp bất cứ lúc nào.

\chenanh[0.65]{10-chon-lop.png}{Màn hình chọn lớp}

\textbf{Bước 3 --- Tham gia lớp trên Google Classroom (nếu giáo viên đã kết
nối).} Sau khi chọn lớp, màn hình hiện mã lớp Classroom kèm hướng dẫn ba bước.
Học sinh mở Google Classroom, bấm dấu cộng, nhập mã lớp là vào được lớp thật
của giáo viên. Bước này không bắt buộc --- có thể bấm \textit{Bỏ qua} để vào
dùng ngay.

\chenanh[0.65]{11-tham-gia-classroom.png}{Màn hình hướng dẫn tham gia lớp Google Classroom kèm mã lớp}

\textbf{Bước 4 --- Tạo đề.} Có hai cách:

\begin{itemize}[leftmargin=1.2cm, itemsep=3pt]
  \item \textbf{Cách 1 --- gõ yêu cầu bằng tiếng Việt} vào ô hội thoại, ví dụ:
        \textit{``Tạo cho em đề kiểm tra 15 phút chương 1 lớp 10''} hoặc
        \textit{``Cho em đề kiểm tra giữa kì 1 môn Toán lớp 10''}. Nên nói rõ
        \textbf{lớp}, \textbf{phạm vi} (chương mấy, giữa kì hay cuối kì) và
        \textbf{loại bài} (kiểm tra thường xuyên hay định kì).
  \item \textbf{Cách 2 --- dùng nút ``Tạo đề nhanh''}, chọn sẵn lớp, chương và
        loại bài trong các ô chọn. Cách này không cần diễn đạt, chạy nhanh hơn
        và luôn cho kết quả đúng ý --- nên dùng khi cách 1 hiểu sai yêu cầu.
\end{itemize}

\noindent\textit{(Màn hình hội thoại này đã trình bày ở Hình~\ref{hinh:01-giao-dien-hoc-sinh.png}.)}

\chenanh[0.65]{13-tao-de-nhanh.png}{Biểu mẫu ``Tạo đề nhanh'' --- chọn lớp, chương, loại bài}

\textbf{Bước 5 --- Nhận đề.} Sau khoảng nửa phút đến một phút (hệ thống phải
biên dịch tệp \LaTeX\ ra PDF), màn hình hiện đường dẫn tải đề. Học sinh có thể
tải PDF về in ra làm giấy, hoặc bấm \textit{Làm bài online} để làm ngay trên
web.

\noindent\textit{(Tệp đề và tệp lời giải dạng PDF xem ở Hình~\ref{hinh:07-de-va-loi-giai-pdf.png}.)}

\textbf{Bước 6 --- Làm bài và xem kết quả.} Trang làm bài hiện lần lượt các câu
trắc nghiệm nhiều lựa chọn, đúng/sai và trả lời ngắn. Làm xong bấm \textit{Nộp
bài}, hệ thống chấm ngay và hiện điểm kèm đáp án đúng của từng câu. Phần tự
luận học sinh làm ra giấy nộp cho giáo viên, nên điểm chấm tự động tối đa là 7
trên thang 10 --- màn hình có ghi rõ điều này.

\noindent\textit{(Trang làm bài và màn hình kết quả xem ở Hình~\ref{hinh:15-lam-bai-nop-bai.png}.)}

\textbf{Bước 7 --- Xem lời giải và luyện lại.} Dưới mỗi đề trong khung hội
thoại có nút \textit{``Lời giải của đề này''} để tải bản lời giải chi tiết đúng
của đề đó. Ở trang kết quả có nút \textit{``Làm đề khác cùng cấu trúc''}: bấm
vào sẽ nhận ngay một đề mới cùng dạng, cùng mức độ, chỉ khác số liệu --- dùng
để luyện lại phần vừa làm sai.

\chenanh[0.65]{16-loi-giai-va-lam-de-khac.png}{Nút tải lời giải của từng đề và nút ``Làm đề khác cùng cấu trúc''}

\textbf{Bước 8 --- Hỏi lại thầy/cô AI câu chưa hiểu.} Ngay dưới mỗi câu ở trang
kết quả có nút \textit{``Hỏi thầy/cô AI về câu này''}. Bấm vào, gõ chỗ mình
vướng (ví dụ: \textit{``em chưa hiểu tại sao lại ra bước này''}), máy sẽ giảng
lại từng bước.

Ba điều học sinh cần biết:

\begin{itemize}[leftmargin=1.2cm, itemsep=2pt]
  \item \textbf{Phải nộp bài rồi mới hỏi được.} Câu chưa làm thì nút bị mờ.
        Tự nghĩ trước thì lúc nghe giảng mới vào đầu.
  \item \textbf{Bên dưới lời giảng luôn có khung xanh ``Lời giải chuẩn của
        thầy/cô''}, gồm đủ đề bài, đáp án và lời giải. Nếu hai phần nói khác
        nhau thì \textbf{tin khung xanh}, vì đó mới là bản do chương trình sinh
        ra. Và báo lại cho thầy cô biết.
  \item \textbf{Mỗi ngày hai mươi lượt hỏi.} Hết lượt thì vẫn đọc được lời giải
        chuẩn, chỉ không có phần giảng lại.
\end{itemize}

Nếu đã rời trang kết quả, học sinh vẫn hỏi lại được: về Chat AI, bấm nút
\textit{``Hỏi lại đề cũ''} ở dưới ô nhập, hệ thống hiện lại đề gần nhất chia
theo Phần I--IV, bấm vào số câu nào là giảng câu đó.

\chenanh[0.65]{23-nut-hoi-lai-de-cu.png}{Nút ``Hỏi lại đề cũ'' trong Chat AI: chọn câu theo từng phần của đề}

\subsubsection{Hướng dẫn dành cho giáo viên}

\textbf{Bước 1 --- Đăng nhập.} Giáo viên đăng nhập bằng chính tài khoản Google
đang dùng cho Google Classroom, để hai bên khớp nhau.

\textbf{Bước 2 --- Kết nối Google Classroom (làm một lần duy nhất).} Vào địa
chỉ \texttt{/gv/classroom/connect}, chọn tài khoản Google và bấm đồng ý cấp
quyền. Hệ thống xin sáu quyền: đọc danh sách lớp, đọc danh sách học sinh, đọc
email học sinh, đăng bài tập, đăng tài liệu, và tải tệp lên Drive của chính
giáo viên. Sau khi đồng ý, màn hình báo kết nối thành công và liệt kê các lớp
đọc được.

\chenanh[0.65]{17-gv-ket-noi-classroom.png}{Màn hình cấp quyền và kết quả kết nối Google Classroom}

\textbf{Bước 3 --- Đồng bộ danh sách lớp.} Bấm đồng bộ để hệ thống lấy về danh
sách lớp và học sinh. Từ đây, mã lớp sẽ tự hiện ra cho học sinh ở Bước 3 phần
trên.

\textbf{Bước 4 --- Lấy đề kiểm tra.} Giáo viên tạo đề giống như học sinh (Bước
4 phần trên), chọn loại bài \textit{kiểm tra định kì} để có đề đủ cấu trúc 12
câu trắc nghiệm, 2 câu đúng/sai, 4 câu trả lời ngắn, 3 câu tự luận. Tạo nhiều
lần sẽ được nhiều mã đề cùng cấu trúc, khác số liệu --- in ra dùng kiểm tra
trực tiếp. Bản lời giải tải riêng bằng nút \textit{``Lời giải của đề này''}.

\textbf{Bước 5 --- Đề, đáp án và điểm về Classroom.} Khi học sinh làm bài
xong, hệ thống tự đăng đề, đáp án và điểm của em đó về lớp Classroom tương
ứng. Học sinh mở Classroom quen thuộc là thấy bài của mình.

\chenanh[0.65]{18-classroom-hoc-sinh.png}{Đề, đáp án và điểm hiện trong tài khoản Google Classroom của học sinh}

\textbf{Bước 6 --- Xem thống kê lớp.} Vào \texttt{/gv/thong-ke} để xem điểm
trung bình của lớp, điểm từng học sinh và những chương, bài học sinh sai nhiều
nhất --- căn cứ để điều chỉnh việc dạy.

\noindent\textit{(Trang thống kê xem ở Hình~\ref{hinh:02-trang-thong-ke-giao-vien.png}.)}

\textbf{Bước 7 --- Đọc nhật kí trợ giảng.} Vào \texttt{/gv/gia-su} để xem học
sinh đã hỏi những gì và máy trả lời ra sao. Bảng có hai cột đặt cạnh nhau:
\textit{AI trả lời} và \textit{Đáp án Python} --- đối chiếu là biết ngay máy có
nói lệch không. Cột cuối cho phép nâng thêm lượt hỏi cho từng em trong ngày.

Trang này còn có nút \textbf{``Kiểm tra kết nối n8n''}. Bấm một cái, hệ thống
gửi sang công cụ điều phối một mã ngẫu nhiên và bắt máy đọc lại; mã quay về
nghĩa là câu lệnh tới nơi, trợ giảng hoạt động đúng. Nếu không quay về, hệ
thống nói thẳng ra chỗ cần sửa. Phép thử này không tốn lượt hỏi của học sinh
nào. Tôi làm nút này sau khi mất mấy ngày dò một lỗi không nhìn thấy được
(lần 19, Bảng 7).

\subsubsection{Chuyển sang tiếng Anh}

\textbf{Học sinh.} Bấm nút \textit{English} ở góc phải trên của trang (cạnh tên tài
khoản); muốn quay lại bấm \textit{Tiếng Việt}. Giao diện, đề bài, lời giải và lời
giảng của trợ giảng đều chuyển sang tiếng Anh. Lưu ý: chỉ những đề tạo \textit{sau}
khi hệ thống có tính năng này mới có bản tiếng Anh; đề cũ, hoặc đề thuộc chương
chưa dịch, vẫn hiện tiếng Việt --- khi đó em bấm \textit{``Làm đề khác''} hoặc tạo đề
mới. Ở đề tiếng Anh, số thập phân viết bằng dấu chấm (\texttt{3.5}).

\textbf{Giáo viên.} Ở biểu mẫu \textit{Ra đề}, tick ô \textit{``Tạo kèm đề tiếng Anh
tương ứng''} nếu cần cả hai thứ tiếng. Trang \textit{Đề đã tạo} có thêm hàng liên
kết \textit{Tiếng Anh} (PDF đề, PDF lời giải, mã nguồn \LaTeX). Đề tiếng Anh và đề
tiếng Việt cùng câu, cùng số liệu nên có thể phát song song cho hai nhóm học sinh
mà vẫn dùng chung một đáp án. Ở khu làm việc của giáo viên, nút ngôn ngữ nằm ngay
dưới thanh menu để không che các mục trên thanh.

\subsubsection{Những lỗi thường gặp và cách khắc phục}

Trong quá trình chạy thử tôi ghi lại mười bốn tình huống người dùng hay gặp.
Bảng dưới đây nêu sáu tình huống gặp nhiều nhất; bảng đầy đủ nằm trong sổ tay kĩ
thuật của hệ thống, có đường dẫn ở cuối mục này.

\begin{center}
\begin{longtable}{|p{4.2cm}|p{4.0cm}|p{5.2cm}|}
\hline
\textbf{Hiện tượng} & \textbf{Nguyên nhân} & \textbf{Cách khắc phục} \\ \hline
\endhead
Gõ yêu cầu nhưng hệ thống hỏi lại ``Em học lớp mấy?'' &
Câu yêu cầu chưa nêu rõ lớp &
Nêu rõ lớp 10, 11 hay 12. Hệ thống cố ý hỏi lại thay vì tự đoán, để tránh ra
nhầm đề \\ \hline

Báo lỗi, không ra được đề &
Câu yêu cầu diễn đạt khó hiểu, hoặc dịch vụ xử lí ngôn ngữ đang bận &
Bấm nút \textit{``Tạo đề nhanh''} hiện ngay dưới thông báo lỗi --- đường đi này
không qua trí tuệ nhân tạo nên luôn chạy được \\ \hline

Trong đề có dòng ``[THIẾU CÂU HỎI]'' &
Ngân hàng chưa có đủ câu ở đúng ô ma trận đó (thường là câu trả lời ngắn mức
vận dụng cao) &
Vẫn dùng được đề, chỉ cần thay câu đó. Đây cũng là chỉ dấu cho giáo viên biết
cần bổ sung hàm sinh nào tiếp theo \\ \hline

Điểm tối đa chỉ tới 7 &
Đúng thiết kế: 3 điểm tự luận làm ra giấy, chức năng chấm tự luận tự động đang
được xây dựng &
Giáo viên chấm phần tự luận và cộng vào. Màn hình kết quả đã ghi rõ để học sinh
không hiểu nhầm \\ \hline

Học sinh không thấy bài trong Classroom &
Học sinh chưa tham gia lớp bằng mã lớp; hoặc giáo viên chưa kết nối Classroom &
Học sinh vào Classroom nhập mã lớp. Giáo viên kiểm tra lại
\texttt{/gv/classroom/connect} \\ \hline

Nút ``Hỏi thầy/cô AI'' bị mờ, không bấm được &
Câu đó học sinh chưa nộp bài, hoặc là câu tự luận chưa có lời giải mẫu &
Làm và nộp bài câu đó trước. Đây là quy định có chủ ý, không phải lỗi \\ \hline
\end{longtable}
\end{center}
\textit{Bảng 8. Sáu tình huống người dùng hay gặp nhất (bảng đầy đủ mười bốn tình huống xem tại sổ tay kĩ thuật)}

\linkqr{https://github.com/lantuan/NganHangDe/tree/main/docs}{Sổ tay kĩ thuật: bảng lỗi đầy đủ, hướng dẫn chi tiết cho giáo viên và học sinh, mô tả kiến trúc hệ thống}

\ghichu{Khi cho đồng nghiệp dùng thử, chị ghi thêm những câu họ hay hỏi vào
bảng đầy đủ trong sổ tay kĩ thuật --- không cần đưa hết vào báo cáo.}

\bigskip
